\documentclass[12pt]{article}
\usepackage[utf8]{inputenc}
\usepackage[T1]{fontenc}
\usepackage{hyperref}
\title{Tanguat Language Resources for K-12 Curriculum Development}
\author{Compiled by Claude Sonnet 5 (claude-sonnet-5) for Andr\'as Kornai}
\date{}
\begin{document}
\maketitle

\section{Introduction}
Tanguat (ISO 639-3 tbs) is a Ramu language of the northern coast of Madang Province, Papua New Guinea, with about 740 speakers (2003, per Wikipedia). Joshua Project says many people also use Tok Pisin. It should not be confused with Tangu (ISO tgu), a separate language of the Bogia District. Joshua Project's description of Tanguat as Austronesian conflicts with the Ramu classification and is probably wrong.

\section{Primary Educational Resources}
No primers, textbooks or school materials were found in search results.

\section{Linguistic and Reference Materials}
\begin{itemize}
\item Wikipedia, Tanguat language: \url{https://en.wikipedia.org/wiki/Tanguat_language}
\item New Guinea World, Ramu River languages (family overview): \url{https://newguineaworld.linguistik.uzh.ch/families/keram-and-ramu-rivers/ramu-river/start}
\end{itemize}

\section{Religious and Community Resources}
\begin{itemize}
\item Joshua Project, Tanguat in Papua New Guinea: \url{https://joshuaproject.net/index.php/people_groups/15255}
\end{itemize}

\section{Digital and Audio Resources}
No audio, video or corpus resources were found in search results.

\section{Curriculum Development Context}
Search turned up no Tanguat teaching materials, orthography or school programme. Tok Pisin and English are the likely school languages in the area. Any native-language curriculum would need new documentation work with the community.
\end{document}